% !TeX root = ../manual.tex
% ===========================================================================
%  Preguntas frecuentes
% ===========================================================================

\section{Preguntas frecuentes}
\label{sec:faq}

\subsection*{¿Olvidé mi contraseña, cómo la recupero?}
\addcontentsline{toc}{subsection}{¿Olvidé mi contraseña, cómo la recupero?}

La aplicación no tiene un flujo de autoservicio para restablecer la
contraseña. Contacta al administrador del club para que verifique tu
identidad y gestione el cambio.

\subsection*{Inicié sesión pero no veo ninguna reserva}

Revisa el filtro activo en \boton{Mis Reservas}: si está en
\boton{Confirmadas}, \boton{Finalizadas} o \boton{Canceladas} y tu reserva
está en otro estado, no aparecerá. Pulsa \boton{Todas} para ver el
historial completo (\cref{sec:usuario:mis-reservas}).

\subsection*{No veo ningún bloque libre en una cancha}

Puede deberse a alguna de estas causas:

\begin{itemize}
  \item Todos los bloques del día ya están ocupados por otras reservas.
  \item Hay un bloqueo de mantenimiento vigente sobre ese rango
        (\cref{sec:admin:mantenimiento}); se ve en color naranja.
  \item La cancha está fuera de su horario de atención para esa hora.
\end{itemize}

Prueba con otro día, u otra cancha del mismo deporte.

\subsection*{Intenté reservar un bloque y me lo rechazó}

Si el mensaje indica que el bloque ya no está disponible, es porque otra
persona lo reservó un instante antes que tú. El sistema garantiza que un
mismo bloque nunca quede reservado dos veces, así que solo una de las dos
solicitudes se acepta. Elige otro bloque u otra cancha.

\subsection*{¿Cuántas reservas puedo tener a la vez?}

Como máximo el tope de reservas activas que el club tenga configurado —3
de forma predeterminada, pero el administrador puede cambiarlo
(\cref{sec:admin:configuracion})—. Si ya llegaste al tope vigente,
cancela primero alguna de las que ya tienes (\cref{sec:usuario:cancelar}).

\subsection*{Quiero cancelar una reserva y no aparece el botón Cancelar}

Solo se pueden cancelar reservas en estado \textsc{Confirmada}. Si el
bloque horario de la reserva ya pasó, queda en estado \textsc{Finalizada}
y deja de poder cancelarse.

\subsection*{Mi cuenta dice «inactiva» y no puedo iniciar sesión}

Un administrador desactivó tu cuenta. Contacta al club para que la
reactive desde Gestión de usuarios (\cref{sec:admin:usuarios}). Tus
reservas anteriores se conservan mientras tanto.

\subsection*{Cancelé una cancha por error, ¿se pierden sus reservas?}

No. Inactivar una cancha solo la retira de Disponibilidad para reservas
nuevas; las reservas ya confirmadas y el historial no se modifican.
Reactívala con el mismo icono para volver a ofrecerla
(\cref{sec:admin:canchas}).

\subsection*{¿Cómo doy de alta a un nuevo administrador?}

Desde la aplicación no es posible: toda cuenta creada con \boton{Crear
cuenta} nace como usuario final (\cref{sec:acceso:crear-cuenta}). Ascender
una cuenta a administrador es una operación que realiza el equipo técnico
del club directamente sobre la base de datos.

\subsection*{Como usuario final, intenté ver Reportes, Canchas, Usuarios o Configuración}

Esas secciones son exclusivas del rol administrador. Si accedes con una
cuenta de usuario final, el sistema rechaza la operación; no forman parte
del menú que ves al iniciar sesión.

\subsection*{¿El administrador también puede reservar desde Disponibilidad?}

No. El administrador tiene su propia sección \boton{Disponibilidad}
(\cref{sec:admin:disponibilidad}), con los mismos bloques que ve un
usuario final, pero es de solo consulta: sirve para resolver dudas o
reclamos, no para crear reservas. Esa acción sigue siendo exclusiva del
usuario final.
